\section{INTRODUCTION}
\label{sec:introduction}

\subsection{Background and Motivation}

Communication and sensing are often needed within the same system. A mobile
terminal needs a data link and knowledge of its position, while a network
operator needs to carry traffic and monitor the infrastructure that supports
it. Integrated sensing and communication (ISAC) brings these functions together
through shared signals, hardware, or infrastructure
\cite{Zhang2026ISACEvolution}. This interest extends to future mobile networks,
with ISAC included among the usage scenarios in the International Mobile
Telecommunications framework for 2030 and beyond
\cite{ITUR_M2160_2023}.

Optical integrated sensing and communication (O-ISAC) applies this idea to
optical links and photonic systems. A fiber can carry communication and sensing
signals along the same route, while a visible light source can support data
transmission and positioning \cite{OISAC_SCR00013,OISAC_SCR00196}.
Figure~\ref{fig:oisac_platform_overview} illustrates representative ways of
combining the two functions. The role of optics differs across these systems.
In fiber and optical wireless links, light propagates through the sensing
medium or reaches the receiver or target. In photonics-enabled wireless
systems, optical components can generate the signal while propagation to the
target takes place in the millimeter-wave or terahertz band
\cite{OISAC_SCR00036,OISAC_SCR00083}.

\begin{figure*}[!t]
\centering
\includegraphics[width=\textwidth]{figures/fig_introduction_oisac_platforms.png}
\caption{Representative O-ISAC configurations in fiber, free-space optical,
visible-light, and photonics-enabled terahertz systems. Each panel shows a
shared resource and the communication and sensing functions it supports.
The design choices below the panels can affect both functions, with their
roles depending on the implementation. These examples illustrate physical
configurations rather than a performance ranking or an exhaustive taxonomy.}
\label{fig:oisac_platform_overview}
\end{figure*}

Sharing a resource can also link the performance of the two functions.
In one fiber experiment, raising the sensing-probe power while holding
communication launch power fixed improved the sensing signal-to-noise ratio
but increased the communication bit error rate \cite{OISAC_SCR00057}.
Joint operation can also allow sensing to assist communication. In a
visible-light prototype, position estimates guided beamforming toward the
receiver and improved the communication data rate at the tested locations
\cite{OISAC_SCR00196}. These examples show why the shared resource and the way
it is used matter. Joint design must account for both the benefit sought and
its effect on the other function.

The physical setting determines how these relationships can be used.
A vibration measurement along a fiber and a position estimate in a room
answer different sensing questions. Optical wireless links also depend on
propagation and alignment, while fiber systems have their own transmission
and sensing constraints \cite{OISAC_SCR00916,OISAC_SCR01085,Mohsan2026OISACReview}.
A researcher selecting an O-ISAC approach therefore needs more than a list
of achievable rates or sensing errors. The useful questions are which design
choice produced the result, how both functions responded, and whether the
conditions match the intended application. This connects the reported
performance to decisions about how to build and evaluate a system.

\subsection{Related Surveys}
\label{sec:related_surveys}

Existing surveys explain O-ISAC from several complementary perspectives.
Architecture-oriented works discuss how communication and sensing can be
combined, while studies of analytical limits and estimation explain the
performance available under particular models
\cite{Wen2024OISACArchitectures,Su2025OISACTheoryPractice,
Khorasgani2026GeneralOpticalISACSurvey}. Other reviews examine visible-light
systems, fiber networks, or photonic waveform generation in greater detail.
Table~\ref{tab:prior_surveys} summarizes their main topics and places the
present survey alongside them.

\begin{table*}[!t]
\caption{Main topics of representative related surveys and the present review}
\label{tab:prior_surveys}
\centering
\begingroup
\footnotesize
\setlength{\tabcolsep}{4pt}
\renewcommand{\arraystretch}{1.08}
\begin{tabularx}{\textwidth}{@{}%
>{\raggedright\arraybackslash}p{0.22\textwidth}
>{\raggedright\arraybackslash}p{0.17\textwidth}
>{\raggedright\arraybackslash}X@{}}
\toprule
Survey focus & Representative works & Main topics \\
\midrule
Architecture and system design &
\cite{Wen2024OISACArchitectures,Su2025OISACTheoryPractice,
Chu2025ROISACOverview,Wang2026CrossLayerISACoF} &
Optical integration architectures, joint design, retroreflective links,
and cross-layer operation in fiber networks. \\
\addlinespace[2pt]
Analytical limits and estimation &
\cite{Khorasgani2026GeneralOpticalISACSurvey,Rode2025PolarizationISAC} &
Communication and sensing models, performance bounds, estimation methods,
and polarization sensing in coherent fiber links. \\
\addlinespace[2pt]
Evolution and deployment &
\cite{Zhang2026ISACEvolution,Liu2025OpticalNetworkISACProspect,
Cao2026SubmarineFiberISAC,Zhang2026OpticalTransmissionNetworks,
Ip2026DeployedTelecomCableISAC} &
The development of ISAC and its standardization, optical network
architectures, and terrestrial and submarine deployment. \\
\addlinespace[2pt]
Physical media and optical networks &
\cite{Liang2024LiSACReview,Zhang2023OpticalNetworkISACFrontier,
He2025FiberISACAdvances,Liu2026OpticalFiberISACReview,
Lu2026ConvergedFiberISACReview,Wang2026FiberNetworkISAC} &
Visible-light integration, fiber sensing mechanisms, and communication
and sensing over optical networks. \\
\addlinespace[2pt]
Waveforms and photonic hardware &
\cite{Yang2025DSCMISACReview,He2026PhotonicFiberISAC,
Yu2023PhotonAssistedTHz,Bai2025MicrowavePhotonicsISAC,
Wang2025MicrowavePhotonicsISAC,Lyu2026PhotonicTHzWaveforms} &
Waveform design, multiplexing, photonic signal generation,
and receiver processing in fiber and wireless systems. \\
\addlinespace[2pt]
Fundamentals and applications &
\cite{Mohsan2026OISACReview} &
Optical platforms, experimental demonstrations, waveform and hardware
design, applications, and research challenges. \\
\midrule
This survey & --- &
Links between physical platforms, shared design choices, and communication
and sensing results, followed by validation needs and research priorities. \\
\bottomrule
\end{tabularx}
\endgroup
\end{table*}

Recent surveys also connect several parts of the field. Zhang et al.
\cite{Zhang2026ISACEvolution} examine the evolution of ISAC across spectrum,
network architecture, and sensing modalities, together with security and
standardization. Mohsan et al. \cite{Mohsan2026OISACReview} review optical
platforms and experimental demonstrations alongside waveform design,
applications, and open challenges. These broad treatments provide a basis for
understanding the available technologies. Our organizing question concerns
the link between a design choice and its communication and sensing outcomes.
We use this link to explain what the reported results mean for a particular
system and which conclusions can inform another design.

\subsection{Scope and Contributions}

This survey examines 206 unique O-ISAC studies first made publicly available
from 1 January 2020 through 22 June 2026. It covers guided optical systems,
optical wireless links, and photonics-enabled wireless systems, including
hybrid configurations. Across these settings, we examine how communication
and sensing are connected through a physical resource, a signal, or a joint
design objective. The review follows three questions. What is shared between
the two functions? How does the shared design affect their performance?
What further evaluation is needed to use the resulting insight in another
system?

The survey makes three contributions.
\begin{itemize}
\item We organize optical platforms by their physical signal paths and the
resources that connect communication and sensing. This distinguishes the
integration mechanism from the hardware or application label.

\item We examine how shared design choices affect communication and sensing
results. We interpret numerical findings within their measurement and
operating conditions and identify the tradeoffs or mutual benefits reported
by the studies.

\item We relate these findings to the way systems have been validated.
This supports research priorities stated in terms of the design to test,
the comparison to make, and the communication and sensing outcomes to
measure.
\end{itemize}

Section~\ref{sec:foundations_comparison} introduces the physical and performance
concepts used in the review, and Section~\ref{sec:review_methods} describes the
review methods. Sections~\ref{sec:optical_modality_families}
through~\ref{sec:technologies_applications_6g} examine platforms, joint design,
validation, and applications. Section~\ref{sec:discussion_roadmap} develops the
research priorities, and Section~\ref{sec:conclusion} concludes the article.

\FloatBarrier
